\section{Theoretical Background}
In this section, we provide the foundational theories and architectures that underpin the proposed Cross-Modality Semantic Alignment Transformer (CMSAT). We first review the fundamental components of the Transformer architecture, followed by an overview of the Zero-Shot Learning (ZSL) framework, and finally discuss the integration of Large Language Models (LLMs) for industrial fault semantics.

The Transformer architecture, originally introduced for machine translation, has revolutionized fields like computer vision and signal processing through its capacity for modeling long-range dependencies and complex structural patterns in sequential data. Its core, the Multi-Head Self-Attention (MHSA) mechanism, allows the model to dynamically weight input sequence segments by evaluating mutual correlations across the entire sequence length. This bypasses the sequential processing limits of recurrence and enables direct global context modeling, which is essential for capturing periodic components in mechanical vibration signals.

Given input embeddings representing temporal or spectral segments, self-attention computes Queries ($Q$), Keys ($K$), and Values ($V$) via linear transformations with weight matrices $W^Q$, $W^K$, and $W^V$. Attention scores derive from the scaled dot product of $Q$ and $K$, followed by a softmax operation to normalize those scores into attention weights. The output is a weighted sum of $V$:
\begin{equation}
\text{Attention}(Q, K, V) = \text{softmax}\left(\frac{QK^T}{\sqrt{d_k}}\right)V
\end{equation}
where $d_k$ is the key dimension. This lets the model focus on relevant temporal or spectral features in machinery signals, effectively capturing subtle fault signatures often masked by measurement noise or overlapping mechanical components.

MHSA executes this process in parallel across $h$ independent heads, each focusing on different representation subspaces of the input sequence. This allows the model to simultaneously attend to diverse information at various positions, significantly enhancing the richness and robustness of the learned features. The heads' outputs are concatenated and projected to form the final representation:
\begin{equation}
\text{MultiHead}(Q, K, V) = \text{Concat}(\text{head}_1, \dots, \text{head}_h)W^O
\end{equation}
where each $\text{head}_i = \text{Attention}(QW_i^Q, KW_i^K, VW_i^V)$.

The Transformer also includes several critical structural components that ensure performance stability:
\begin{itemize}
    \item \textbf{Positional Encoding:} Since the architecture lacks inherent order knowledge, sinusoidal or learnable encodings are added to inputs to inject relative or absolute position information, preserving the temporal structure of the raw signal.
    \item \textbf{Feed-Forward Networks (FFN):} Each layer contains a position-wise fully-connected FFN consisting of two linear transformations with a non-linear activation (ReLU or GELU) between them, enabling complex feature mapping at each position.
    \item \textbf{Layer Normalization and Residual Connections:} To stabilize training, prevent internal covariate shift, and mitigate vanishing gradients in deep stacks, residual connections are used around each sub-layer, followed by Layer Normalization (LN).
\end{itemize}
These elements collectively enable the extraction of robust hierarchical features from non-stationary rotating machinery signals, making it an ideal investigative backbone for cross-modality alignment tasks \cite{transformer_orig, vaswani2017attention}.

\subsection{Zero-Shot Learning Framework}
Zero-Shot Learning (ZSL) aims to recognize objects from classes unseen during the model's training phase. In industrial fault diagnosis, this paradigm is vital as labeled training data for every failure mode under all possible operating conditions is often impractical, expensive, or inherently dangerous to collect. ZSL addresses this by leveraging auxiliary semantic information to transfer knowledge from known, well-documented scenarios to novel, unseen fault types.

The standard ZSL framework partitions label space into seen classes $\mathcal{Y}_s$, which are available during training, and unseen classes $\mathcal{Y}_u$, which are only encountered during inference. By definition, $\mathcal{Y}_s \cap \mathcal{Y}_u = \emptyset$. Generalization depends on an auxiliary semantic space $\mathcal{S}$ where each class $y$ is represented by a semantic prototype $a_y$ encoding physical attributes, functional characteristics, or textual descriptions.

ZSL learns a mapping between signal feature space $\mathcal{X}$ and semantic space $\mathcal{S}$, often via a compatibility function $f(x, a_y; \theta)$. During training, parameters $\theta$ are optimized to maximize compatibility scores for correct seen classes. For a test sample $x$ belonging to an unseen category, classification selects the class whose prototype $a_y$ has the highest compatibility score within the unseen candidate pool:
\begin{equation}
\hat{y} = \arg\max_{y \in \mathcal{Y}_u} f(x, a_y; \theta)
\end{equation}
Traditional ZSL often suffers from "domain shift," where seen-class mappings fail to generalize accurately to unseen class distributions. A key bottleneck is the quality of semantic descriptors; manual attribute labeling is labor-intensive, requires expert knowledge, and is often prone to subjective bias \cite{palatucci2009zero, lampert2009learning}. We address this limitation by using Large Language Models (LLMs) to generate high-fidelity, standardized semantic prototypes for mechanical faults \cite{xian2018zero}.

\subsection{Large Language Models in Industrial Semantics}
Large Language Models, like the GPT or BERT series, are trained on vast, heterogeneous datasets including technical manuals, engineering reports, materials science, and scientific publications. This extensive pre-training captures technical nuances and engineering knowledge missing from standard word-level embeddings like Word2Vec. LLMs possess internal representations linking technical concepts to physical phenomena.

For fault diagnosis, LLMs bridge the semantic gap between abstract categorical labels (e.g., "Inner Race Fault") and their actual physical manifestations in sensor data. Prompting strategies like Chain-of-Thought (CoT) or structured templates allow LLMs to generate detailed descriptions of how faults affect machinery behavior, covering frequency modulation, energy distribution, and impact pulse patterns. For example, a carefully crafted prompt might request the distinctive vibration characteristics, harmonics, and temporal impulses of a high-speed bearing fault.

These descriptions are transformed into dense semantic embeddings. Two representation levels are typically considered in this context:
\begin{itemize}
    \item \textbf{Word-Level Embeddings:} These represent individual words but often miss the context-dependent meaning of complex technical phrases and engineering jargon common in mechanical engineering.
    \item \textbf{Sentence-Level and Document-Level Semantic Vectors:} Using LLM hidden states or specialized sentence-transformers, we extract a fixed-dimensional vector representing the global semantic essence of a description, capturing conceptual interactions and overall framework.
\end{itemize}
These LLM-generated vectors serve as anchors in our cross-modality framework. Unlike manual attributes, these embeddings contain rich, high-dimensional info about physical nature and dynamics, enabling more discriminative and generalizable mappings between vibration signals and technical definitions \cite{brown2020language, devlin2018bert, industrial_llm_ref}. Aligning signal features with LLM-enhanced semantics allows our model to achieve superior zero-shot performance across varying conditions and novel fault modes. This approach moves beyond simple label matching to a deeper, knowledge-driven alignment of physical signals and technical semantics.
